\uuid{h9kX}
\titre{Matrice/Application linéaire/Noyau/Image/Dimension}
\auteur{hanani ; mzari}
\organisation{exo7}
\licence{CC BY-NC-SA 4.0}
\datecreate{2020-01-20}
\niveau{L1}
\module{Algèbre}
\chapitre{Matrice}
\sousChapitre{Propriétés élémentaires, généralités}
\difficulte{3}
\theme{Calcul matriciel, Matrice, Application linéaire, Noyau, Image, Dimension}
\exo7qcmid{lille-2.matrices.015}

\contenu{
\begin{question}
\qclassification{108.01}
\qtags{motcle=Matrice/Application linéaire/Noyau/Image/Dimension}
Soit $ M_2(\Rr)$ l'ensemble des matrices carrées d'ordre $2$ à coefficients réels et $f$ l'application  définie par :
$$\begin{array}{rccc}f:&M_2(\Rr)&\to& \Rr\\
& M = \left(\begin{array}{rc} a&b\\
c&d\\  \end{array}\right) &\to &\mbox{tr}(M) = a+d, \end{array}$$
le réel $\mbox{tr}(M)$ est appelée la trace de $M$. Quelles sont les assertions vraies ?
\begin{answers}
\good{$f$ est une application linéaire}.
\good{$\dim \ker f = 3$}.
\bad{$\dim \Im f = 2$}.
\good{$\Im f = \Rr$}.
\end{answers}
\begin{explanations} On vérifie que $f$ est une application linéaire et que
$$\ker f= \left\{\left(\begin{array}{rc}
a&b\\c&-a\\ \end{array}\right)\; ;\; a,b,c \in \Rr \right\}.$$
Donc la famille $\left\{\left(\begin{array}{rc}
1&0\\ 0&-1\\ \end{array}\right), \;  \left(\begin{array}{rc}
0&1\\ 0&0\\ \end{array}\right), \; \left(\begin{array}{rc}
0&0\\ 1&0\\ \end{array}\right) \right\}$ est une base de $\ker f$ et $\dim \ker f=3$. Or, d'après le théorème du rang, $\dim \Im f=1=\dim \Rr$ et comme $\Im f$ est un sous-espace vectoriel de $\Rr$, donc $\Im f=\Rr$.
\end{explanations}
\end{question}
}
